% LVS Theory paper - Nexearch P01
\documentclass[11pt]{article}
\usepackage[utf8]{inputenc}
\usepackage{amsmath,amssymb,amsthm}
\usepackage{graphicx}
\usepackage{hyperref}
\usepackage{booktabs}
\usepackage[margin=1in]{geometry}
\usepackage{braket}

\title{Latent Vacuum Stationarity: Reality as Fixed Points of the Quantum Vacuum}
\author{Fabien}
\date{\today}

\begin{document}
\maketitle

\begin{abstract}
We propose that observable physical reality corresponds to the stationary configurations
(fixed points) of the quantum vacuum landscape. This meta-interpretive framework,
\textbf{Latent Vacuum Stationarity} (LVS), unifies the fine-tuning problem, the measurement
problem, and the problem of time through a single principle: \emph{stability is the criterion
of existence}. We show compatibility with 9 empirical tests (6 confirmed, 3 supported) and
demonstrate that the Shaposhnikov-Wetterich fixed-point condition naturally predicts
$m_H \approx 126$ GeV, within 1 GeV of the measured value.
\end{abstract}

\section{Introduction}
% See LVS_Paper_v2.md sections 1-3

\section{The LVS Framework}
The core postulate:
\begin{equation}
\text{Observable reality} \longleftrightarrow \text{Fixed points of } \mathcal{V}
\end{equation}
where $\mathcal{V}$ is the quantum vacuum configuration space.

The Wheeler-DeWitt equation establishes the atemporal substrate:
\begin{equation}
\hat{H}\ket{\Psi} = 0
\end{equation}

Time emerges via Page-Wootters mechanism:
\begin{equation}
\ket{\psi(t)}_S = C\braket{t|\Psi}
\end{equation}

\section{Higgs Mass from Stationarity}
% See LVS_Formalization.ipynb
The fixed-point condition at the Planck scale:
\begin{align}
\lambda(M_{Pl}) &= 0 \\
\beta_\lambda(M_{Pl}) &= 0
\end{align}
Running the SM RG equations down to the electroweak scale yields:
\begin{equation}
m_H \approx 126 \text{ GeV} \quad (\text{measured: } 125.25 \pm 0.17 \text{ GeV})
\end{equation}

\section{Empirical Validation}
% See LVS_Validation.ipynb - 9 tests

\section{Discussion}
% See discussion.md

\section{Conclusion}

\bibliographystyle{plain}
\bibliography{refs}
\end{document}
